% abstract-env.tex — CUMCM 国赛摘要环境

\begin{abstract}
\noindent

% === 摘要内容由 P5a 生成后填入此区域 ===
% 550-700 字，漏斗式结构：
%   钩子(问题本质+为什么难) →
%   核心洞察(用什么视角/工具破解) →
%   模型骨架(逻辑关系不是清单) →
%   关键结果(有对比基准的数字) →
%   结论延伸价值

% 注意：
%   - 不能出现"本文""我们"之外的人称
%   - 不能出现无对比基准的裸数字
%   - 不能罗列模型名而无逻辑关系

\textbf{关键词：}关键词1；关键词2；关键词3；关键词4

\end{abstract}
